The hypothesis that our universe is a computer simulation is widely discussed but rarely stated in a form that can be tested. We review the philosophical arguments for it, its antecedents in digital physics, and every empirical signature proposed in the literature that we could identify. These signatures are a spacetime lattice, finite spatial resolution, render-on-demand quantum measurement, bounded computational resources, and information-physical effects. For each one we state the auxiliary assumptions it needs, the observations that constrain it, and the non-simulation physics that predicts the same signature. We argue that the unrestricted hypothesis cannot be refuted by any observation. Only conjunctions with explicit assumptions about the simulator's architecture and resources, and about whether it adapts to being probed, can be tested. For every such sub-hypothesis we found, a detection would establish new physics, not simulation, and a null result bounds only a class of simulator designs. The review also contains four analyses of our own. First, applying current Lorentz-violation bounds to the lattice dispersion relations of Beane, Davoudi and Savage raises the lower limit on the inverse lattice spacing of their unimproved-Wilson scenario from $\sim10^{11}$~GeV to $\gtrsim10^{16}$~GeV. We also find that photon decay into $e^+e^-$ is kinematically allowed far below the lattice scale, contrary to the original analysis. Second, we find no published search for the cubic-symmetric cosmic-ray anisotropy that scenario predicts, and we specify one that uses public data. Third, we show that the success criteria of the proposed render-on-demand experiments contradict standard quantum mechanics. Fourth, we reproduce the SARS-CoV-2 entropy values underlying the proposed ``second law of infodynamics'' and show that the reported decrease is accounted for by the known mutational bias of the virus. Every factual claim in the paper is backed by a public ledger that records its source, location and supporting quotation.
